\documentclass[11pt]{article}

\usepackage[utf8]{inputenc}
\usepackage[T1]{fontenc}
\usepackage{geometry}
\usepackage{amsmath}
\usepackage{amssymb}
\usepackage{booktabs}
\usepackage{hyperref}
\usepackage{xurl}
\usepackage{xcolor}
\usepackage{listings}
\usepackage{enumitem}
\usepackage{parskip}
\usepackage{graphicx}
\graphicspath{{../images/}}

\geometry{margin=1in}

\lstset{
  language=Java,
  basicstyle=\ttfamily\small,
  keywordstyle=\color{blue},
  commentstyle=\color{gray},
  stringstyle=\color{teal},
  showstringspaces=false,
  columns=fullflexible,
  frame=single,
  framerule=0.2pt,
  breaklines=true
}

\setlist[itemize]{topsep=0.3em,itemsep=0.2em}
\setlist[enumerate]{topsep=0.3em,itemsep=0.2em}

% Simple per-section source line.
\newcommand{\notesources}[1]{\par\noindent{\small\color{gray}\textit{Sources:} \sloppy #1\par}}

% Unit-wise source strings for this subject (used by slides/notes).
% Path is relative to the lecture `latex/` folder (the compilation CWD).
% OOPS with Java (CSEG1044) — unit-wise sources
%
% These are short, student-facing reference strings intended to be used with:
%   - \slidesources{...}  (in beamer slides)
%   - \notesources{...}   (in lecture notes)
%
% Style inspiration: other subjects in My-Subjects/ use semicolon-separated sources
% and include an "accessed" date for web documentation.

\newcommand{\OOPSUnitISources}{Schildt, \textit{Java: A Beginner's Guide} (10e), McGraw-Hill (Java fundamentals + OOP basics).; Oracle Java Tutorials: Getting Started + Language Basics + Classes and Objects (accessed \today).; Java SE API docs (e.g., \texttt{String}, \texttt{Scanner}) (accessed \today).}

\newcommand{\OOPSUnitIISources}{Schildt, \textit{Java: The Complete Reference} (12e), McGraw-Hill (classes, inheritance, packages, interfaces).; Bloch, \textit{Effective Java} (3e), Addison-Wesley, 2018 (design choices; interfaces vs abstract classes).; Oracle Java Tutorials: Classes and Objects + Interfaces + Packages (accessed \today).; Java SE API docs (e.g., \texttt{Comparable}, \texttt{Runnable}) (accessed \today).}

\newcommand{\OOPSUnitIIISources}{Schildt, \textit{Java: The Complete Reference} (12e) (nested/inner classes, exceptions, threads, I/O).; Oracle Java Tutorials: Nested Classes + Exceptions + Concurrency + Basic I/O (accessed \today).; Java SE API docs (e.g., \texttt{Thread}, \texttt{AutoCloseable}) (accessed \today).}

\newcommand{\OOPSUnitIVSources}{Schildt, \textit{Java: The Complete Reference} (12e) (generics, lambdas, Swing, JDBC).; Bloch, \textit{Effective Java} (3e) (generics + lambdas).; Oracle Java Tutorials: Generics + Lambda Expressions + Swing + JDBC Basics (accessed \today).; Java SE API docs (e.g., \texttt{java.util.function}, \texttt{javax.swing}, \texttt{java.sql}) (accessed \today).}

\newcommand{\OOPSUnitVSources}{Schildt, \textit{Java: The Complete Reference} (12e) (Collections Framework + wrapper classes).; Oracle Java docs: Collections Framework overview (accessed \today).; Java SE API docs (\texttt{java.util} collections; wrapper classes in \texttt{java.lang}) (accessed \today).}

\newcommand{\OOPSUnitVISources}{Git SCM: \textit{Pro Git} (2e) (version control workflow), accessed \today.; GitHub Docs (repositories, issues, pull requests), accessed \today.; Oracle Java Tutorials: Swing + JDBC (accessed \today).; JUnit 5 User Guide (accessed \today).; Apache Maven Project: Getting Started (accessed \today).; Gradle User Manual: Getting Started (accessed \today).; UML class diagrams: UML 2.x overview/spec (accessed \today).}

\newcommand{\OOPSMiscWhyInterfacesSources}{Schildt, \textit{Java: The Complete Reference} (12e) (interfaces).; Bloch, \textit{Effective Java} (3e) (prefer interfaces where appropriate).; Oracle Java Tutorials: Interfaces (accessed \today).; Java SE API docs (e.g., \texttt{Comparable}, \texttt{Runnable}) (accessed \today).}



\title{Object Oriented Programming (CSEG1044)\\Unit 02 -- Lecture 03 Notes\\Access Modifiers + Packages + Interface Intro}
\begin{document}
\maketitle
\tableofcontents

\section{Lecture Snapshot}
\begin{itemize}
  \item \textbf{Unit focus:} Classes, Inheritance, Packages, and Interfaces
  \item \textbf{Main new ideas:} Access control (what each means), Packages: why and how, Interface intro: contract view
  \item \textbf{Course thread:} Use Access Modifiers + Packages + Interface Intro to strengthen the domain model, APIs, and access rules inside Campus Resource Manager.
\end{itemize}
\notesources{\OOPSUnitIISources}

\section{Lecture Overview}
As your codebase grows, you need two things:
\begin{itemize}
  \item \textbf{control visibility} of classes/fields/methods (to prevent misuse),
  \item \textbf{organize code} into packages (so files and names stay manageable).
\end{itemize}
We also introduce the \textbf{interface} idea as a contract to reduce coupling.
\notesources{\OOPSUnitIISources}

\section{Prerequisite Bridge}
Before reading the new material in depth, do a fast recall of the older ideas listed below.
\begin{itemize}
  \item \textbf{Encapsulation}: Access modifiers are the concrete mechanism for the encapsulation idea. Main reference: \texttt{Unit 02 / Lecture 01 slides/notes}.
  \item \textbf{Class organization}: Packages are easier to understand once class members and static utilities are familiar. Main reference: \texttt{Unit 02 / Lecture 02 slides/notes}.
\end{itemize}
This lecture only revisits those ideas briefly so that the main explanation can stay focused on the new concept sequence.
\notesources{\OOPSUnitIISources}

\section{Learning Outcomes}
After this lecture, you should be able to:
\begin{itemize}
  \item Choose correct access modifiers to protect a class design.
  \item Organize code using packages and imports.
  \item Explain interfaces as contracts and write a basic interface example.
\end{itemize}

\section{Suggested Reading Order}
\begin{enumerate}
  \item Read the \textbf{Prerequisite Bridge} first and confirm each recalled idea.
  \item Study the central explanation in this order: \textit{Access control (what each means)} $\rightarrow$ \textit{Packages: why and how} $\rightarrow$ \textit{Interface intro: contract view}.
  \item Trace the worked example line by line before checking short answers.
  \item Attempt the practice section without looking at the solution immediately.
  \item Finish with the exit questions and further-study suggestions to confirm retention.
\end{enumerate}
\notesources{\OOPSUnitIISources}

\section{Key Terms (Quick)}
\begin{itemize}
  \item \textbf{Package}: namespace that groups related classes (also maps to folder structure).
  \item \textbf{Package-private (default)}: visible only inside the same package.
  \item \textbf{Protected}: visible in same package and in subclasses (even in other packages).
  \item \textbf{Public API}: the classes/methods you intentionally expose for others to use.
  \item \textbf{Coupling}: how strongly one part of code depends on another.
\end{itemize}

\section{Detailed Notes}
\subsection{Access control (what each means)}
Visibility rules (simplified for day-to-day use):
\begin{itemize}
  \item \texttt{private}: only inside the same class.
  \item \textbf{default} (no keyword): inside the same package.
  \item \texttt{protected}: inside the same package, and also in subclasses.
  \item \texttt{public}: everywhere.
\end{itemize}

\paragraph{Rule of thumb.}
\begin{itemize}
  \item Keep fields \texttt{private}.
  \item Expose methods only if necessary.
  \item Use package-private for helper classes that should not be part of the public API.
\end{itemize}

\subsection{Packages: why and how}
Why packages?
\begin{itemize}
  \item Avoid name conflicts (\texttt{Student} in two different domains).
  \item Improve organization and readability.
  \item Control visibility using package-private access.
\end{itemize}

\textbf{How to declare a package:} the first non-comment line should be:
\begin{itemize}
  \item \texttt{package com.upes.crm.model;}
\end{itemize}
The folder path should match:
\begin{itemize}
  \item \texttt{com/upes/crm/model/Student.java}
\end{itemize}

\textbf{Imports:}
\begin{itemize}
  \item Use \texttt{import} to avoid writing fully-qualified names.
  \item Example: \texttt{import java.util.List;}
\end{itemize}

\subsection{Interface intro: contract view}
An interface defines \textbf{what operations} a class must provide, without specifying how.
This allows you to program to an interface (a contract), not to a concrete class.

\textbf{Key benefit:} reduced coupling. Your code depends on an abstract contract, so implementations can change without changing the caller.

We will cover interfaces deeply later (multiple files, design decisions, and real use cases).
\notesources{\OOPSUnitIISources}

\section{Running Project Connection}
Use Access Modifiers + Packages + Interface Intro to strengthen the domain model, APIs, and access rules inside Campus Resource Manager.
\begin{itemize}
  \item Use Campus Resource Manager as the running case study, or map the same idea to your own project domain if you are using another example.
  \item Ask where today's idea should live: model, service, DAO, UI, or team workflow.
  \item Use the lecture example as a smaller version of the same design choice you will make in the capstone.
\end{itemize}
\notesources{\OOPSUnitIISources}

\section{Common Mistakes and Confusions}
\begin{itemize}
  \item Using public as a shortcut whenever code fails to compile.
  \item Creating packages with no responsibility boundary.
  \item Treating an interface as just a weaker class.
\end{itemize}
\notesources{\OOPSUnitIISources}

\section{Worked Example: packages + interface (small)}
\subsection*{File 1: \texttt{com/upes/demo/pay/Payable.java}}
\begin{lstlisting}
package com.upes.demo.pay;

public interface Payable {
    double amountDue();
}
\end{lstlisting}

\subsection*{File 2: \texttt{com/upes/demo/pay/Invoice.java}}
\begin{lstlisting}
package com.upes.demo.pay;

public class Invoice implements Payable {
    private double total;

    public Invoice(double total) {
        if (total < 0) throw new IllegalArgumentException("total");
        this.total = total;
    }

    @Override
    public double amountDue() {
        return total;
    }
}
\end{lstlisting}

\subsection*{File 3: \texttt{com/upes/demo/app/Main.java}}
\begin{lstlisting}
package com.upes.demo.app;

import com.upes.demo.pay.Invoice;
import com.upes.demo.pay.Payable;

public class Main {
    public static void main(String[] args) {
        Payable p = new Invoice(2500.0);
        System.out.println(p.amountDue());
    }
}
\end{lstlisting}
The line \texttt{Payable p = new Invoice(...)} is used here only to show the \textbf{contract view}: the caller depends on \texttt{amountDue()}, not on invoice internals. The runtime dispatch explanation comes in Lecture 04.

\section{Demo / Observation Task}
Use this block as a short reproducible demo or observation task.
\begin{itemize}
  \item Reproduce the lecture's main example around \textit{Access control (what each means)} once without looking at the solution first.
  \item Change one input, rule, or design choice in Campus Resource Manager and predict the result before checking it.
  \item Record one success case, one edge case, and one failure case so the idea becomes observable instead of purely theoretical.
\end{itemize}
\notesources{\OOPSUnitIISources}

\section{Practice Check (with brief solutions)}
\begin{enumerate}
  \item Which access modifier should be used for most fields?\; \textbf{Answer:} \texttt{private}.
  \item If a member has no access modifier, where is it visible?\; \textbf{Answer:} within the same package.
  \item Why do we use packages?\; \textbf{Answer:} organization + avoid name conflicts + control visibility.
\end{enumerate}

\section{Self-Study Practice (no solutions)}
\begin{enumerate}
  \item Create package \texttt{com.upes.bank} with classes \texttt{Account} and \texttt{Customer}.
  \item Write an interface \texttt{Printable} with method \texttt{print()} and implement it in two classes.
  \item Make a helper class package-private and observe where it can/can't be used.
\end{enumerate}

\section{Further Study}
\begin{itemize}
  \item Revisit the prerequisite references if any recall bullet still feels weak.
  \item Rework the lecture example with one changed assumption, input, or design choice.
  \item Read the textbook and documentation references listed in the source line for a second explanation of the same topic.
  \item Apply the same idea once inside Campus Resource Manager so that the concept moves from theory to design practice.
\end{itemize}
\notesources{\OOPSUnitIISources}

\section{Exit Questions (with sample answers)}
\textbf{Q1.} What is the difference between \texttt{private} and package-private?\par
\textbf{Sample answer:} \texttt{private} is visible only inside the same class. Package-private (default) is visible to any class in the same package.

\vspace{0.6em}
\textbf{Q2.} Why are interfaces useful?\par
\textbf{Sample answer:} Interfaces define a contract and reduce coupling. Code written against an interface can work with multiple implementations without changing the caller.

\vspace{0.6em}
\textbf{Q3.} Why do packages matter even in a small project?\par
\textbf{Sample answer:} They control visibility, communicate architecture, and reduce naming collisions.

\end{document}
